\section{本章小结}
\label{sec:sgen-summary}

本章从独立均匀的基础随机数出发，回答了章首的问题：怎样把它们加工成给定分布的样本，以及加工结果是什么性质的样本。方法的选择取决于目标分布提供了哪些可以计算的信息。表~\ref{tab:sgen-methods}把各种方法放在一起比较。

\begin{BookInlineTable}
  \centering
  \small
  \begin{tabular}{@{}llll@{}}
    \toprule
    方法 & 需要的信息 & 输出 & 主要成本与失效 \\
    \midrule
    逆变换 & 分位数或可反演的$F$ & 独立样本 & 数值反演误差、尾部截断 \\
    仿射高斯 & 协方差因子 & 独立样本 & 一次分解，奇异时退化 \\
    分步生成 & 成分或条件分布 & 独立样本 & 各步须用不同随机数 \\
    拒绝采样 & $\widetilde{p}$、建议、包络$M$ & 独立样本 & 平均$M/Z$次试验 \\
    MH & $\widetilde{p}$、建议密度 & 相关的链 & 步长取舍、支持连通 \\
    Gibbs & 完全条件分布 & 相关的链 & 强相关时慢、支持形状 \\
    Hamilton & $\nabla U$、光滑正密度 & 相关的链 & 每步$L$次梯度、步长稳定 \\
    \bottomrule
  \end{tabular}
  \caption{采样生成方法的比较。前四种输出服从目标分布的独立样本；后三种输出以目标为不变分布的链，状态相关，能否走遍目标需另行论证。}
  \label{tab:sgen-methods}
\end{BookInlineTable}

直接生成的三种机制都以推前为核心。变量变换把易抽变量送到目标，正确性由原像规则或可逆映射的换元公式验证；逆变换对一切实值分布成立，分位数可以解析计算时最简便，需要数值反演时，引理~\ref{lem:sgen-inversion-error}说明逐点误差$\delta$怎样把分布函数夹在$F(x\pm\delta)$之间，前提是误差界对所有$u$都成立。高斯样本由独立标准坐标经仿射变换得到，协方差奇异时输出落在低维子空间上，没有密度但分布正确。

分步生成依靠全概率公式与链式分解，关键条件是各步使用互不重叠的随机数；输入独立并不意味着同一样本的坐标独立。拒绝采样只需未归一化密度与有效包络，接受样本精确服从目标，代价是平均$M/Z$次试验，在高维中常因包络过松而无法使用。只要每次都用一组新的基础随机数从头生成，以上方法都给出独立同分布的目标样本。

链式更新换了一种思路：放弃一次造出完整样本，转而构造以目标为不变分布的转移。不变性是一个代数条件，细致平衡是验证它的便利充分条件。Metropolis--Hastings用接受概率$\min\{1,\widetilde{p}(y)q(x\mid y)/(\widetilde{p}(x)q(y\mid x))\}$平衡双向概率流，其中建议比修正了建议本身的偏向，归一化常数约去；Gibbs采样按完全条件分布逐坐标更新，选坐标的概率不能依赖当前状态；Hamilton采样借助辅助动量和保持体积、可逆的蛙跳格式走出远距离的提议，再用能量差修正离散误差。三者都只证明了不变性。

不变性不保证从任意初值到达目标：停在原地的转移、无法越过空隙的建议、被支持形状挡住的坐标更新都是反例。即使从目标出发，链的状态也相互相关。

核验生成结果时，要分别检查输出分布或不变性的证明、链的可达性、停止规则和数值近似带来的偏差。直方图和样本矩能发现错误，但看不到从未访问的区域，也不能代替证明。

链的长期行为，即在什么条件下分布收敛、时间平均收敛以及收敛有多快，由卷九讨论。下一章转向用这些样本计算积分：独立样本平均的误差、近似生成造成的偏差、链初值的影响，以及相关样本平均的方差，都将在那里分别评价。
